\documentclass[11pt, a4paper]{../problems_template}

\begin{document}

\begin{titlepage}
  \begin{center}
	\Huge{Математический анализ}
  \end{center}
\end{titlepage}

\begin{exercise}
Пусть $X$ - непустое множество. На множестве его подмножеств $2^{X}$ опеределена такая операция "звездочка" $A \rightarrow A^{\star}$, что $A^{\star} \cup A = X$ и $(A^{\star})^{\star} = A$ для всякого $A \subset X$. Обязана ли операция "звездочка" совпадать с операцией дополнения, то есть $A^{\star} = X \setminus A$?
\end{exercise}

\end{document} 
